\documentclass[10pt,a4paper]{amsart}
\usepackage[margin=19mm]{geometry}
\usepackage{amsmath,amssymb,mathtools}
\usepackage[scr=rsfs,cal=euler]{mathalfa}
\usepackage{xfrac}
\usepackage[unicode,hidelinks,bookmarks=false]{hyperref}
\newcommand{\ie}{i.e.\ }
\newcommand{\cf}{cf.\ }
\newcommand{\resp}{respectively\ }
\newcommand{\Subsection}{\subsection}
\providecommand{\nobreakdash}{\nobreak\hbox{-}}
\setlength{\emergencystretch}{2em}
\multlinegap0pt
\usepackage{bm}
\usepackage{bbm}
\usepackage{dsfont}
\usepackage{xspace}
\usepackage{cancel}
\usepackage{mathtools}
\usepackage{amsthm}
\makeatletter
\@ifundefined{proposition}{\newtheorem{proposition}{Proposition}}{}
\makeatother
\title[First-Order Trajectory Matching: Fast Ensemble Predictions o]{First-Order Trajectory Matching: Fast Ensemble Predictions of Chaotic, Turbulent, Stochastic Systems}
\author{Shreya Jha, Timo Schorlepp, Nicholas Geissler, Jules Berman, Benjamin Peherstorfer}
\date{Source: arXiv:2606.11138v1 (CC BY 4.0)}
\begin{document}
\maketitle

\noindent\emph{Source and licence.} First-Order Trajectory Matching: Fast Ensemble Predictions of Chaotic, Turbulent, Stochastic Systems. Version arXiv:2606.11138v1 (CC BY 4.0), distributed under the Creative Commons Attribution 4.0 International licence, as recorded in the arXiv licence metadata for this exact version and shown on the abstract page https://arxiv.org/abs/2606.11138v1. Full rights notice: licenses/arxiv-2606.11138-CC-BY-4.0.txt.\par\smallskip

\noindent\textit{Editorial scope.} Counted unit: the Proposition whose source label is \texttt{None} in the article (source0.tex lines 1306--1328, source label \texttt{None}), delivered together with the article's own complete original proof of it (source0.tex lines 1330--1395, one \texttt{proof} environment of 66 lines). No mathematics is written, repaired or re-derived: statement and proof are the article's own text line for line. Required context retained uncounted: none. Every definition, display and label of those lines is retained unchanged. Omitted from this package and not redistributed: 1--1305, 1329--1329, 1396--3088. Nothing inside a retained excerpt is omitted or altered: every retained body is delivered exactly as the article's lines are written, so no omission receipt of any kind accompanies this package. Citations: the retained excerpts carry no citation command at all, so no bibliography entry of the article is transcribed and no reference is redistributed. Reference adaptation: none was needed and none was made; every reference inside the delivered unit resolves inside it, so no unresolved reference or citation is produced. Printed slips kept literal: no printed word, symbol, hypothesis or number of the counted statement or of its proof was altered. Layout and class: the article's own class is \texttt{article} with packages including caption, blindtext, jcustom, lastpage, inputenc, fontenc, booktabs, hyperref, amsfonts, amsmath, amsthm, amssymb, nicefrac, microtype; the wrapper is the lane's pinned \texttt{amsart} editorial wrapper at 10pt on A4, which restates the theorem-like environments the retained text needs and adds the article's own macro definitions that the retained text needs (none), with extra wrapper packages (bm, bbm, dsfont, xspace, cancel, mathtools). The delivered numbering is the unit's own. Assets: no retained span contains a figure environment, a graphics-inclusion command, a PDF-page inclusion, a table or a TikZ picture; no image file is copied into this package, the package declares no asset dependency and no image file is redistributed with it. Licence: version arXiv:2606.11138v1 of this article is distributed under the Creative Commons Attribution 4.0 International licence, as recorded in the arXiv licence metadata for this exact version and shown on the abstract page https://arxiv.org/abs/2606.11138v1. A full rights notice is delivered at licenses/arxiv-2606.11138-CC-BY-4.0.txt. Characters: every retained excerpt is pure ASCII, so the delivered engine, XeLaTeX with its default Latin Modern text font, reports no missing character.
\begin{proposition}
Assume gradient drift $b(t,x) = -\nabla U(t,x)$, and that the diffusion matrix is isotropic, $\Sigma(t) = \sigma(t)^2 I_d$.
Then:
\begin{enumerate}
    \item the current velocity \(v\) is a gradient field,
    \begin{equation*}
        v(t,x)
        = -\nabla U(t,x) - \frac{\sigma(t)^2}{2}\nabla \log \rho(t,x)
        = -\nabla\!\left(U(t,x) + \frac{\sigma(t)^2}{2}\log \rho(t,x)\right)\,.
    \end{equation*}
    \item \(v\) is the unique minimizer of the kinetic energy functional
    \begin{equation*}
        \mathcal{K}(u)
        :=
        \int_0^T \int_{\mathbb{R}^d} \rho(t,x)\,\|u(t,x)\|^2_2 \,\,\mathrm{d} x\,\,\mathrm{d} t,
    \end{equation*}
    among all velocity fields \(u\) satisfying the continuity equation
    \begin{equation*}
        \partial_t \rho + \nabla \cdot (\rho u) = 0.
    \end{equation*}
\end{enumerate}
Hence, in this setting, the probability current velocity \(v\) coincides with the minimal kinetic energy velocity field associated with the prescribed marginal path \((\rho(t))_{t\in[0,T]}\).
\end{proposition}

\begin{proof}
Under the assumptions \(b=-\nabla U\) and \(\Sigma=\sigma^2 I_d\), we compute
\begin{align*}
    v(t,x)
    &= b(t,x) - \frac12 \Sigma(t)\nabla \log \rho(t,x) = -\nabla U(t,x) - \frac{\sigma(t)^2}{2}\nabla \log \rho(t,x) \\
    &= -\nabla\!\left(U(t,x) + \frac{\sigma(t)^2}{2}\log \rho(t,x)\right).
\end{align*}
Thus \(v\) is indeed a gradient field.

Next, \(v\) is admissible for the marginal-matching problem because it satisfies
\begin{equation*}
    \partial_t \rho + \nabla\cdot(\rho v)=0.
\end{equation*}
We now show that among all admissible vector fields, there is at most one admissible gradient field. Suppose
\[
u_1 = \nabla \phi_1,
\qquad
u_2 = \nabla \phi_2
\]
both satisfy
\begin{equation*}
    \partial_t \rho + \nabla\cdot(\rho u_i)=0,
    \qquad i=1,2.
\end{equation*}
Subtracting the two equations gives
\begin{equation*}
    \nabla\cdot\bigl(\rho (u_1-u_2)\bigr)=0.
\end{equation*}
Let
\[
w := u_1-u_2 = \nabla \psi,
\qquad
\psi := \phi_1-\phi_2.
\]
Then
\begin{equation*}
    \nabla\cdot(\rho \nabla \psi)=0.
\end{equation*}
Assuming standard decay at infinity (or, alternatively, no-flux boundary conditions on a bounded domain), we multiply by \(\psi\) and integrate:
\begin{align*}
    0
    = \int_{\mathbb{R}^d} \psi \,\nabla\cdot(\rho \nabla \psi)\,\,\mathrm{d} x
    = - \int_{\mathbb{R}^d} \rho\, \lVert \nabla \psi \rVert_2^2 \,\,\mathrm{d} x.
\end{align*}
Hence
\[
\nabla \psi = 0
\qquad \text{a.e. on } \{\rho>0\},
\]
and therefore
\[
u_1=u_2
\qquad \text{a.e. with respect to } \rho(t,x)\,\,\mathrm{d} x\,\,\mathrm{d} t.
\]
So there is at most one admissible gradient field.

Finally, the minimizer of the kinetic energy functional
\[
\mathcal{K}(u)=\int_0^T \int_{\mathbb{R}^d} \rho\, \|u\|^2_2 \,\,\mathrm{d} x\,\,\mathrm{d} t
\]
subject to
\[
\partial_t\rho+\nabla\cdot(\rho u)=0
\]
is known to be a gradient field. Since \(v\) is both admissible and a gradient field, and such a field is unique, it follows that \(v\) must coincide with the unique minimizer. Therefore \(v\) is exactly the minimal kinetic energy velocity field.
\end{proof}

\end{document}
